\chapter{Key Visualizations and Representative Trajectories}
\label{chap:visualizations}

While the aggregate metrics describe sweeping structural tendencies, deeper phenomenological insights are found within individual agent trajectories. To prevent the exhaustion common to exhaustive diagnostic logs, this section explicitly highlights the most defining, statistically representative runs from both the benchmark and ablation matrices. Together, they form a cohesive narrative of survival policies under non-stationary weather.

\section{The Baseline PPO Trajectory in Fixed Climates}
Under perfectly deterministic, stationary environments, the Proximal Policy Optimization agent rapidly discovers the optimal mathematical ceiling. 

\begin{figure}[H]
    \centering
    \includegraphics[width=0.9\textwidth]{figures/per_run/001_fertilization_ppo_adaptive_fixed_weather_years_1000_seed_0__diagnostics_panel.png}
    \caption{Diagnostics Panel for the defining PPO run under Fixed Weather. The episode lengths quickly saturate to the physical maximum, while training rewards strictly adhere to a monotonically increasing limit curve. The primary metric showcases an aggressively optimized, nearly monotonic climb to convergence.}
    \label{fig:representative_fixed_ppo}
\end{figure}

\section{The Drought Crash Phenomenon}
The fundamental shift occurs when the same agents are subjected to deeply non-stationary randomized weather patterns simulating 1,000 years of chaotic climate variation. The structural survival policies inevitably encounter severe drought cycles that violently break traditional exploitation policies.

\begin{figure}[H]
    \centering
    \includegraphics[width=0.9\textwidth]{figures/per_run/003_p1_ppo_adaptive_random_weather_seed0_ent0__episode_length_vs_global_step.png}
    \caption{The 'Drought Crash' Phenomenon visible in Episode Length progression. Despite an otherwise climbing optimization curve, injected multi-year historical droughts physically break the carrying capacity of the environment, causing sudden precipitous drops in sequence lengths. Agents relying heavily on zero-entropy exploration are most vulnerable to these shocks.}
    \label{fig:drought_crash}
\end{figure}

\section{Catastrophic Over-Fertilization (Economic Misweighting)}
When economic safety guardrails are artificially reduced to $\texttt{cost\_weight=0.8}$, the agent behavior strictly deviates from profitability towards pure maximal biological yields, effectively operating without regard to the cost of intervention.

\begin{figure}[H]
    \centering
    \includegraphics[width=0.9\textwidth]{figures/per_run/028_p3_ppo_adaptive_random_weather_seed0_costw0_8__primary_metric_vs_global_step.png}
    \caption{Catastrophic Over-Fertilization timeline. Notice the intensely dense activation graph indicating the agent actively 'panic-spraying' inputs. While yielding densely populated crop structures, the economic efficiency fundamentally collapsed under realistic constraints.}
    \label{fig:over_fertilization}
\end{figure}

\section{A2C Volatility under Hierarchical Constraints}
While the Advantage Actor-Critic (A2C) agents nominally navigated the environment, their gradient updates lacked the unclipped constraint bounding of PPO. 

\begin{figure}[H]
    \centering
    \includegraphics[width=0.9\textwidth]{figures/per_run/013_p2_a2c_fixed_weather_seed0_blockpen0__checkpoint_eval_curves.png}
    \caption{A2C Evaluation Checkpoint Curve. Notice the intense parameter oscillation and massive variance bands, indicating volatile learning jumps that periodically regress the agent far below optimal state values prior to rebounding.}
    \label{fig:a2c_volatility}
\end{figure}

\begin{figure}[H]
    \centering
    \includegraphics[width=0.9\textwidth]{figures/per_run/013_p2_a2c_fixed_weather_seed0_blockpen0__crop_decision_timeline.png}
    \caption{The hierarchical crop decision timeline generated by the A2C run mapped physically against internal phase transitions. Despite volatility, the final sequence successfully honors structural sequence laws.}
    \label{fig:a2c_hierarchy}
\end{figure}

\section{Summary Note}
The trajectory mappings validate that aggregate averages do not fundamentally misrepresent underlying training sequences. PPO is strictly superior to A2C internally, and environment-conditional exploration metrics are vital for navigating catastrophic events like randomized drought.
